\section{Eliminating every arithmetic macro}
\label{sec:primitives}

This section specifies the fixed finite graph certificates used by the
direct materialization. All coordinates here are nonnegative integers.
A displayed product or square is subsequently expressed by addition and
multiplication gates; the final gate rows are quadratic.

\subsection{A positive-power graph with 27 auxiliary coordinates}

For positive inputs $n,k$ and proposed output $m$, write
\[
 \mathcal N(a,x,y)=x^2+y^2-1-a^2y^2.
\]
The positive-power graph has the following 27 additional coordinates:
\[
\begin{gathered}
 W,A,T,Z,X,Y,U,V,S,R,B,\ a_0,y_0,b_0,q,d,\\
 b_+,b_-,s_+,s_-,r_+,r_-,x_+,x_-,g,w_n,w_k.
\end{gathered}
\]
Impose the finite equations
\begin{align}
 A&=a_0+2,&Y&=k+y_0,&B&=b_0+2,\label{eq:powbounds}\\
 \mathcal N(A,X,Y)&=0,&
 \mathcal N(A,U,V)&=0,&
 \mathcal N(B,S,R)&=0,\label{eq:pownorms}\\
 B&=1+4Yq,& V&=Y^2(d+1),\label{eq:powdiv}\\
 B+Ub_+&=A+Ub_-,&
 S+Us_+&=X+Us_-,&
 R+4Yr_+&=k+4Yr_-,\label{eq:powcong}\\
 X+Yn+Tx_+&=YA+m+Tx_-,\label{eq:powfinalcong}\\
 2An&=T+n^2+1,& T&=m+g+1,\label{eq:powmodulus}\\
 W&=n+w_n,&W&=k+w_k,&
 \mathcal N(W+1,A,WZ)&=0.\label{eq:powlast}
\end{align}
Existence of these coordinates is equivalent to $m=n^k$.
This is the positive specialization of the Pell--Matiyasevich
power criterion~\cite{pell,carneiro}, with every strict inequality
represented by a natural slack, every divisibility represented by a
quotient, and every possibly signed congruence quotient represented as
a difference of two naturals. The first three norms and congruences
characterize the $k$th Pell solution at parameter $A$. The last norm,
the bounds on $W$, the modulus definition, and the strict output bound
are precisely the remaining power criterion.

The condition $Y\ge k\ge1$ excludes the exceptional zero-index Pell
solution. The modulus $4Y$ is positive. Also \eqref{eq:powlast} gives
\[
 A^2=1+W^3(W+2)Z^2.
\]
Since $A\ge2$ and $W\ge n\ge1$, one has $Z\ge1$ and $A>W\ge n$.
Thus the integer rearrangement in \eqref{eq:powfinalcong} is equivalent
to the original natural-subtraction congruence involving $A-n$.
These side conditions are part of the proof, not optional simplifications.

The optional unrestricted builder retains the zero-exponent and zero-base
branches. The delivered constructions do not need those branches, because
the preceding section and an independently checked call-site induction
establish positive inputs to all 55 power uses.

\subsection{Binomial coefficients by a base digit}

For $n,k\in\N$, put
\[
 q=2^n+1,\qquad L=q^k,\qquad P=(q+1)^n.
\]
The binomial theorem expands $P$ in base $q$, with the coefficient
$\binom nk$ in digit position $k$, because all coefficients are at most
$2^n<q$. In particular, positions above $n$ have digit zero.
For proposed output $c$, add just four natural coordinates $h,r,r_s,c_s$
besides the three power graphs, and impose
\begin{equation}\label{eq:choosegraph}
 P=L(qh+c)+r,\qquad L=r+r_s+1,\qquad q=c+c_s+1.
\end{equation}
The last two equations say $0\le r<L$ and $0\le c<q$. Uniqueness of
quotient, remainder, and base digit proves that \eqref{eq:choosegraph}
holds exactly when $c=\binom nk$.

\subsection{Factorials as a ceiling quotient}

For $n\in\N$, define
\[
 \rho=(n+2)^{n+2},\qquad c=\binom{\rho+n}{n},\qquad P=(\rho+1)^n.
\]
Then
\begin{equation}\label{eq:factorialformula}
 n!=\left\lceil\frac Pc\right\rceil.
\end{equation}
For $n=0,1$ this is immediate. For $n\ge2$,
\[
 \frac Pc=n!\prod_{i=1}^n\frac{\rho+1}{\rho+i}
 =n!\prod_{i=1}^n\left(1-\frac{i-1}{\rho+i}\right).
\]
Using $1-\prod(1-a_i)\le\sum a_i$ for $0\le a_i\le1$, we obtain
\[
 0\le n!-\frac Pc
 <\frac{n!\,n(n-1)}{2(\rho+1)}<1.
\]
The final strict bound follows, for example, from
$n!\,n(n-1)/2<n^{n+2}<(n+2)^{n+2}=\rho$.
This proves \eqref{eq:factorialformula}, including the case where the
quotient is already integral.

For proposed output $f$ and two natural coordinates $r,s$, impose
\[
 P+c=cf+r+1,\qquad c=r+s+1.
\]
Since $c>0$, these say $c(f-1)<P\le cf$, which is equivalent to
$f=\lceil P/c\rceil$. The required powers and binomial graph are those
already specified.

\subsection{Ordinary quadratic rows and one quartic}

Every arithmetic-circuit addition or multiplication is represented by a
new natural coordinate $z$ and the row
\[
 z-x-y=0\qquad\hbox{or}\qquad z-xy=0.
\]
Numerals and shared nodes are treated once at each distinct evaluation.
Signed equality is kept as an equality between two nonnegative
circuits. Each graph obligation is then an ordinary residual of degree
at most two. A finite conjunction of rows $R_\ell=0$ is equivalent over
$\N$ to the one equation $\sum_\ell R_\ell^2=0$.

Before affine elimination, the resulting fixed system has 6,183
coordinates and 5,523 rows. Every residual coefficient has absolute
value at most four. Its expanded squared sum has 17,430 monomials and
coefficient height 1,993. These are measured exact coefficient
dictionaries, rather than an asymptotic bound.

\subsection{Removing nonnegative affine gates}

Suppose a retained gate has the defining row $z-L=0$, where $L$ is an
affine polynomial with nonnegative integer coefficients in earlier
coordinates. Substituting $L$ everywhere and deleting this defining
row and coordinate is a bijection on natural zero sets: the inverse
restores $z=L$. Substitution of an affine polynomial preserves the
quadratic degree bound.

We apply this operation in topological order when the expanded affine
replacement has at most four terms. The support threshold controls
coefficient and expression growth; it is unrelated to correctness.
Exactly 2,569 coordinates and rows disappear. This yields
\[
 6183-2569=3614,\qquad 5523-2569=2954.
\]
The resulting quartic has 15,996 monomials and coefficient height
127,778,545, which is less than $2^{27}$. The increase in coefficient
height is the explicit cost of removing these affine wires.

\begin{table}[h]
\centering
\begin{tabular}{lrr}
\toprule
Measured quantity & Retained wires & Affine elimination\\
\midrule
Outer coordinates & 5 & 5\\
Other witness coordinates & 1,665 & 1,665\\
Gate coordinates & 4,513 & 1,944\\
Power / binomial / factorial graphs & 55 / 16 / 2 & 55 / 16 / 2\\
Quadratic rows & 5,523 & 2,954\\
Expanded monomials & 17,430 & 15,996\\
Maximum residual coefficient & 4 & 4,556\\
Maximum quartic coefficient & 1,993 & 127,778,545\\
\bottomrule
\end{tabular}
\caption{Exactly measured direct materializations of the same certificate.}
\end{table}
